% -----------------------------------------------------------------------
% \tablefillin - nachruestbarer Befehl fuer Luecken IN TABELLEN.
%
% \fillin's eingebaute Unterstreichung (\hrulefill) ist in Tabellenzellen
% redundant zum Zellenrand. Dieser Befehl uebernimmt die Antwortlogik von
% \fillin (\ifprintanswers, \CorrectChoice@Emphasis / \color@begingroup /
% \color@endgroup fuer konsistente Antwortfarbe), laesst aber \hrulefill
% und \answerclearance weg.
%
% WICHTIG: Im Nicht-Antwort-Zweig muss \phantom{#1} verwendet werden,
% NICHT leer lassen - sonst kollabiert die tabularx-Spaltenbreite im
% Schueler-Modus (ohne Antworten).
%
% EINBAU: Dieser Block gehoert DIREKT in schulexam.cls selbst, direkt
% neben die bestehende \fillin-Definition - nicht in einzelnen
% Arbeitsblatt-Dateien wiederholen. Kein \makeatletter/\makeatother
% noetig, wenn er direkt in der .cls-Datei steht (dort ist @ bereits
% ein Letter-Charakter).
% -----------------------------------------------------------------------

\newcommand\tablefillin[1]{%
  \leavevmode
  \ifprintanswers
    \color@begingroup\CorrectChoice@Emphasis #1\color@endgroup
  \else
    \phantom{#1}%
  \fi
}% tablefillin

% Verwendung im Dokument (nach Einbau in schulexam.cls), z.B. in einer
% tabularx-Zelle:
%   \tablefillin{42}
